\clearpage
\setcounter{chapter}{7}
\chaptermark{光全息术}
\pdfbookmark[0]{第七章 光全息术}{chapter7}
\problem{7-1}{
以全息术的眼光重新看待两束平行光的干涉。如图（a）所示，两束相干的平行光束可分别看作物光波 $\widetilde O$ 和参考光波 $\widetilde R$，两者波长相同即其波矢均为 $k$ 值，且方向平行于 $(xz)$ 平面，与纵轴 $z$ 之夹角分别为 $\theta_O$ 和 $\theta_R$。
\begin{enumerate}
\item 试写出 $\widetilde O$ 波与 $\widetilde R$ 波的干涉场在全息干版 $H$ 面上的波前函数 $\widetilde U(x,y)$。
\item 求全息干版上所呈现的干涉条纹的间距 $\Delta x$，当 $\theta_O=\theta_R=1^\circ$ 时，或当 $\theta_O=\theta_R=60^\circ$ 时。设光波长为 $633\,\mathrm{nm}$。
\item 某感光胶片厂生产了一种记录干版，其性能为：感光层厚度 $l$ 约 $8\,\mu\mathrm m$，分辨率 $2000$ 线/$\mathrm{mm}$。问，若选择 $\theta_O=\theta_R=60^\circ$ 来拍摄这干涉场，该记录介质的分辨率是否适用？
\item 经线性冲洗后获得一张全息图。若照明的平行光束 $\widetilde R'$ 沿原记录时的参考光 $\widetilde R$ 方向，斜入射于这张全息图，如图（b）所示，试写出再现光的波前函数 $\widetilde U_H(x,y)$，并从中分析出该衍射波场的主要成分及其特点。
\item 若照明光波 $\widetilde R'$ 正入射于这张全息图，如图（c）所示，试写出再现光的波前函数 $\widetilde U_H(x,y)$，并从中分析出该衍射波场的主要成分及其特点。
\end{enumerate}
\tikzfig{fig_7_1}{7-1}
}
\problem{7-2}{
以全息术的眼光重新看待球面波与平面波的干涉。如图（a）所示，正入射的一束平行光作为一个参考光波 $\widetilde R$，与轴外点源发射的傍轴球面波 $\widetilde O$ 相干叠加于记录介质 $H$ 平面。
\begin{enumerate}
\item 试写出 $\widetilde O$ 波与 $\widetilde R$ 波的干涉场在 $H$ 面上的波前函数 $\widetilde U(x,y)$，设物点位置坐标为 $(x_0,0,z_0)$；波长为 $\lambda$。
\item 分析全息干版 $H$ 上所呈现的干涉条纹的特征。
\item 记录介质经曝光和线性冲洗后成为一张全息图。若用一束平行光 $\widetilde R'$ 正入射于这张全息图，如图（b）所示，试写出再现光的波前函数 $\widetilde U_H(x,y)$，并从中分析出该衍射波场的主要成分及其特点。要求作图示意。
\item 若用一束傍轴球面波 $\widetilde R'$ 照明这张全息图，如图（c）所示，试写出再现光的波前函数 $\widetilde U_H(x,y)$，并从中分析出该衍射波场的主要成分及其特点。设照明点源的纵向距离为 $z'$。
\item 试问，在（4）情况下，再现两列发散球面波是可能的吗？再现两列会聚球面波是可能吗？提示：比较 $z_0,z'$ 之大小。
\item 当用球面波照明全息图时，这球面波就相当于一个透镜，它作用于平面波照明时再现的物光波 $+1$ 级虚像上和共轭波 $-1$ 级实像上，从而导致最终成像的放大或缩小。现以本题（4）和图（c）为例，试给出横向放大率 $V_{+1}$ 和 $V_{-1}$。这里，定义 $V_{+1}=x_{+1}/x_0$，$V_{-1}=x_{-1}/x_0$。
\end{enumerate}
\tikzfig{fig_7_2}{7-2}
}
\problem{7-3}{
改变照明光 $\widetilde R'$ 的波长也可能造成全息成像的缩放。现以习题 7-2 图（a）、（b）为实例。设全息记录时的光波长为 $\lambda$，全息再现时的照明光波长为 $\lambda'$。
\begin{enumerate}
\item 试求 $+1$ 级虚像和 $-1$ 级实像的位置坐标 $(x_{+1},0,z_{+1})$ 和 $(x_{-1},0,z_{-1})$，并进而讨论纵向放大率 $M_{\pm1}=z_{\pm1}/z_0$ 和横向放大率 $V_{\pm1}=x_{\pm1}/x_0$。
\item 若照明光波 $\widetilde R'$ 为一束斜入射的平面波，其与 $z$ 轴倾角为 $\theta'$，光波长为 $\lambda'$，试讨论放大率 $M_{\pm1}$ 和 $V_{\pm1}$ 有何变化。
\end{enumerate}
\tikzfig{fig_7_2}{7-2}
}
\problem{7-4}{
讨论题：若用非单色光记录、用非单色光照明，将究竟出现怎样的图景。现以习题 7-1 图（a）、（b）为实例，设全息记录时 $\widetilde O$ 光和 $\widetilde R$ 光均含两种波长 $\lambda_1$ 和 $\lambda_2$，照明光波 $\widetilde R'$ 也含同样的两种波长。为简单起见，设它们的振幅均相等，即
\[
A_O(\lambda_1)=A_O(\lambda_2)=A_R(\lambda_1)=A_R(\lambda_2)
=A'_R(\lambda_1)=A'_R(\lambda_2)=A.
\]
\begin{enumerate}
\item 试给出拍摄时记录介质平面上的波前函数 $\widetilde U(x,y)$ 和光强分布 $I(x,y)$。
\item 试写出经曝光和线性冲洗后所获得的这张全息图的透过率函数 $\widetilde t_H(x,y)$。
\item 当 $\widetilde R'$ 光照明这张全息图时将生成怎样的波前函数 $\widetilde U_H(x,y)$，分析出与其相对应的衍射场的主要成分；并与单色光记录和照明情形作比较，在目前非单色光条件下出现了怎样复杂的图景？
\end{enumerate}
\tikzfig{fig_7_1}{7-1}
}
